% This file should be replaced with your file with thesis content.
%=========================================================================
% Authors: Michal Bidlo, Bohuslav Křena, Jaroslav Dytrych, Petr Veigend and Adam Herout 2019

% For compilation piecewise (see projekt.tex), it is necessary to uncomment it and change
% \documentclass[../projekt.tex]{subfiles}
% \begin{document}

% ===================================
%           INTRO
% ===================================

\chapter{Introduction}
Audio classification is a crucial task in digital signal processing, involving the categorization of audio signals based on their inherent acoustic characteristics.
With the integration of audio-driven technologies, from voice-activated systems to automated media monitoring, the problem of speech and music classification remains relevant even today.
Automation of speech and music classification is now a crucial building block of real-world applications, including the preprocessing of signals for Automatic Speech Recognition (ASR)~\cite{BENZEGHIBA2007763}, the organization of vast multimedia archives~\cite{samouelian1998speech}, and the optimization of telecommunication bandwidth through content-specific approaches, such as Google's Lyra~\cite{ren2024fewertokenneuralspeechcodec} for speech or Microsoft's DCVC~\cite{jia2025practicalrealtimeneuralvideo} for video.
Because speech and music possess distinct statistical and structural properties, identifying the content type is a vital first step, as subsequent processing and encoding operations rely on different algorithms optimized for each category.

Historically, the demand for lightweight classifiers, combined with limited computational resources and skepticism toward neural networks (NN) during periods often referred to as "AI winters" \cite{CognitechSystems2024}, led researchers to develop speech and music classifiers based primarily on classical machine learning (ML) techniques. 
These traditional methods often relied on manual feature engineering, utilizing descriptors such as Zero-Crossing Rate (ZCR), Spectral Flux, and Pitch Variance. 
Notable examples include the decision tree classifier proposed by Lavner et al.~\cite{Lavner2009-zo}, and systems utilizing Gaussian Mixture Models (GMM) and Support Vector Machines (SVM) by Khonglah et al.~\cite{KHONGLAH201671}. 
Although effective in controlled environments, these classical approaches often struggle with the complexity and noise found in real-world audio.

With the revival of neural networks in the 2010s, driven by increased computational power and the availability of large-scale datasets, research into audio classification was revitalized. Deep learning (DL) has significantly advanced the field by enabling models to automatically extract high-level, meaningful features from raw audio or time-frequency representations, such as Mel-Frequency Cepstral Coefficients (MFCC).
Although modern deep learning models have achieved high accuracy, the search for more efficient, robust, and accessible methods continues, particularly for deployment in resource-constrained environments.

The goal of this thesis is to explore the current state of speech and music classification, investigate various machine learning and deep learning architectures, and develop a novel classification method that improves existing performance. 
The thesis begins with an overview of the acoustic differences between speech and music and a review of the popular ML and DL methods used for their classification. 
The data acquisition process is then detailed, including the creation of a large-scale, diverse dataset designed to improve model generalization. 
As a primary contribution, this work presents a new classification approach and evaluates its performance against established baselines. 
Furthermore, the proposed classifiers have been implemented in a ready-to-use software library, which facilitates their practical application. 
Experimental results demonstrate that the proposed method achieves competitive accuracy and robustness, offering a streamlined solution for modern audio analysis tasks.

% ===================================
%           THEORY
% ===================================

\chapter{Speech and Music Classification}
\label{cap:theory_speech_music}
Automatic classification of speech and music originated as a practical quality-of-life feature in radio broadcasting, allowing listeners to automatically mute advertisements or switch between different types of content~\cite{543290}. 
Today, speech–music classification has evolved into a fundamental component of modern audio processing ecosystems. 
It is used across a wide range of technologies, from voice-activated intelligent assistants that must reliably identify human speech to conserve power and improve command accuracy to music discovery platforms that detect and exclude spoken commentary to provide a seamless listening experience.

The significance of speech–music classification is particularly highlighted in telecommunication systems, where speech and music are processed using different codecs tailored to their unique acoustic characteristics.
This requires precise content identification to guarantee high perceptual quality at low bitrates.
Beyond consumer electronics, classification systems play a critical role in large-scale infrastructures such as automated broadcast monitoring, digital library indexing, and real-time audio surveillance, where the ability to categorize large volumes of audio data efficiently is essential for both operational efficiency and content analysis.

\section{Classical Machine Learning Approaches}
\label{sec:theory_classical}
Early speech and music classification systems relied on classical machine learning models because they were SOTA at the time and were less computationally demanding.
They often utilized the extraction of all possible features, possibly extended by statistics of said features aggregated over a sliding window.
The features needed to be experimentally hand-picked to address hardware constraints or selected algorithmically.

DT proposed by Lavner et al.~\cite{Lavner2009-zo} and GMM and SVM from Khonglah et al.'s paper~\cite{KHONGLAH201671} were re-implemented and used as reference in this work.

\section{Neural Network Approaches}
\label{sec:theory_nn}
With the advancement of hardware and the resurgence of neural networks in the 2010s, the field shifted toward neural-based models for the task of speech/music classification.

The approaches typically utilize MFCCs with additional features fed to CNNs, Recurrent Neural Networks (RNN), their combination, CRNN, and others. 
The range of possibilities with different models and features has exploded.

The building block of today's largely used Large Language Models, Transformers~\cite{vaswani2023attentionneed}, did not overlook this field either.

A CRNN created by Choi et al.~\cite{choi2016convolutionalrecurrentneuralnetworks} and Gong et al.'s Audio Spectrogram Transformer (AST)~\cite{gong2021astaudiospectrogramtransformer}, were picked as reference models for this thesis.

\section{State-of-the-Art in Audio Classification}
\label{sec:theory_sota}
Current state-of-the-art systems leverage large-scale datasets and pre-trained audio models to improve robustness and generalization.
Prominent architectures differ fundamentally in their input representations: self-supervised approaches, such as wav2vec~\cite {baevski2020wav2vec20frameworkselfsupervised}, directly process raw audio waveforms, whereas weakly supervised models, like Whisper~\cite{radford2022robustspeechrecognitionlargescale}, utilize log-mel spectrograms as input.
Although these deep models dominate benchmark performance, efficiency and deployability remain key challenges, motivating ongoing research into lightweight and hybrid classification approaches.

% ===================================
%           DATASET
% ===================================

\chapter{Dataset}
\label{cap:dataset}

The data set aggregates publicly available audio sources (see \cref{sec:dataset_sources}), all released under the MIT license or compatible open licenses.

Weak labels were converted into strong, time-aligned annotations by automatically detecting active and silent regions. Class-specific activity detection was applied:

\begin{itemize}
    \item \textbf{Speech:} Silero VAD~\footnote{\url{https://github.com/snakers4/silero-vad}} for the detection of speech activity.
    \item \textbf{Music:} PyDub~\footnote{\url{https://github.com/jiaaro/pydub}} amplitude-based silence detection.
\end{itemize}

\section{Sources}
\label{sec:dataset_sources}
The data set combines speech, music, and mixed-content repositories:

\begin{itemize}
    \item \textbf{Speech:}
    \begin{itemize}
        \item English--Arabic Speech Translation (English subset)~\footnote{\url{https://hf.co/datasets/ammagra/english-arabic-speech-translation}}
        \item ATCO2 Corpus~\footnote{\url{https://hf.co/datasets/Jzuluaga/atco2_corpus_1h}}
        \item NOIZEUS~\footnote{\url{https://ecs.utdallas.edu/loizou/speech/noizeus/}}
    \end{itemize}

    \item \textbf{Music:}
    \begin{itemize}
        \item MusicGen Fine-tuning (Jazz, Electronic, Classical)~\footnote{\url{https://hf.co/datasets/PlutoG99001/}}
        \item Lewtun Music Genres~\footnote{\url{https://hf.co/datasets/lewtun/music_genres}}
        \item Country Music Genres~\footnote{\url{https://hf.co/datasets/ylacombe/music_genres_Country}}
        \item POPMUSIC1981~\footnote{\url{https://hf.co/datasets/memepottaboah/POPMUSIC1981}}
        \item CCMusic Acapella~\footnote{\url{https://hf.co/datasets/ccmusic-database/acapella}}
    \end{itemize}

    \item \textbf{Mixed:}
    \begin{itemize}
        \item MUSAN~\footnote{\url{https://www.openslr.org/17/}}
    \end{itemize}
\end{itemize}

\section{Statistics}
\label{sec:dataset_statistics}
All data were randomly shuffled and split into disjoint subsets:

\begin{itemize}
    \item \textbf{Training:} 80\%
    \item \textbf{Validation:} 10\%
    \item \textbf{Test:} 10\%
\end{itemize}



% ===================================
%           PIPELINE
% ===================================

\chapter{Pipeline}
\label{cap:pipeline}

% ===================================
%           MODELS
% ===================================
\chapter{Models}
\label{cap:models}

% ===================================
%           EXPERIMENTS
% ===================================
\chapter{Experiments}
\label{cap:experiments}

% ===================================
%           CONCLUSION
% ===================================
\chapter{Conclusion}
\label{cap:conclusion}


%=========================================================================

% For compilation piecewise (see projekt.tex), it is necessary to uncomment it
% \end{document}
